% At the very start of section1.tex, before subsection 1.1,
% replace the current first sentence with:

This paper argues that physical models can fail in two distinct ways:
predictively, when their equations give wrong numerical answers within
their domain, and ontologically, when the kind of entity they
presuppose does not exist at the relevant scale.
The classical calculation of the radiative collapse of the hydrogen
atom is our primary case study---a derivation that is mathematically
exact, internally consistent, and empirically false---from which we
derive a general diagnostic applicable to any physical model.
The same structure appears in the ultraviolet catastrophe and in the
quantum measurement problem, as we note in Section~2.
We then argue that Aristotle's distinction between potentiality and
actuality, recovered by Heisenberg in 1958, provides a more adequate
ontological framework for the quantum atom than either classical
particle realism or strict instrumentalism.

% Section .............................................................................................................
\section{The Classical Collapse: An Explicit Calculation}
\label{section:collapse}

The failure of the planetary models of the atom is a well known and extensively discussed example. We present the calculation in full because the philosophical argument that follows depends on it: the failure is not a rough estimate but an exact consequence of three internally consistent and independently verified physical laws.

% Subsection ..........................................................................................................
\subsection{The Three Ingredients and Their Sources}
\label{subsection:ingredients}

\textbf{Centripetal acceleration (Newton, 1687)}

\medskip

Newton's \emph{Principia Mathematica} \cite{newton_1687_principia} establishes that for uniform circular motion of radius $r$ and speed $v$, differentiating the position vector $\vec{r}(t) = r(\cos\omega t,\,\sin\omega t)$ twice with $\omega = v/r$ gives
\begin{equation}
  \ddot{\vec{r}} = -\omega^2\,\vec{r} \quad\Longrightarrow\quad \boxed{a = \frac{v^2}{r}} \; ,
  \label{eq:centripetal}
\end{equation}
directed toward the centre. No approximation is involved.

\medskip

\textbf{Coulomb's law (Coulomb, 1785)}

\medskip

Measured by torsion balance in 1785 \cite{coulomb_1785_electricity}, the force arising between two point charges $q_1$ and $q_2$ separated by distance $r$ is
\begin{equation}
  \boxed{F = \frac{1}{4\pi\epsilon_0}\frac{q_1\,q_2}{r^2}} \; ,
  \label{eq:coulomb}
\end{equation}
attractive for opposite charges. The constant $\epsilon_0$ would be added by Maxwell's 1865; the inverse-square dependence is Coulomb's direct experimental result.

\medskip

\textbf{Larmor radiation formula (Larmor, 1897)}

\medskip

From Maxwell's equations, the electric field radiated by an accelerated charge falls as $1/r$, so the Poynting flux integrated over a distant sphere yields a finite radiated power \cite{larmor_1897_radiation}:
\begin{equation}
  \boxed{P = \frac{e^2\,a^2}{6\pi\epsilon_0\,c^3}} \; .
  \label{eq:larmor}
\end{equation}
The dependence on $a^2$ follows because the radiated field is linear in $a$, and power is quadratic in the field amplitude.

% Subsection ..........................................................................................................
\subsection{The Differential Equation for \texorpdfstring{$r(t)$}{r(t)}}
\label{subsection:ode}

Setting Coulomb's force equal to the centripetal requirement for the electron (mass $m_e$, charge $e$) one obtains:
\begin{equation}
  \frac{e^2}{4\pi\epsilon_0\,r^2} = \frac{m_e\,v^2}{r} \implies v^2 = \frac{e^2}{4\pi\epsilon_0\,m_e\,r} \; .
  \label{eq:virial-v}
\end{equation}

The total mechanical energy is then:
\begin{equation}
  E = \frac{1}{2}m_e v^2 - \frac{e^2}{4\pi\epsilon_0\,r} = -\frac{e^2}{8\pi\epsilon_0\,r} \; ,
  \label{eq:energy}
\end{equation}
where the virial theorem ensures $E_{\mathrm{kin}} = -E/2$.

\medskip

As the electron loses energy by radiation, $E$ decreases and $r$ decreases monotonically. Substituting $a = e^2/(4\pi\epsilon_0 m_e r^2)$ into~(\ref{eq:larmor}):
\begin{equation}
  P = \frac{e^6}{12\pi^3\epsilon_0^3\,m_e^2\,c^3\,r^4}.
  \label{eq:power-r}
\end{equation}

Equating total energy $\dot{E} = -P$ and differentiating~(\ref{eq:energy}) one obtains:
\begin{equation}
  \frac{e^2}{8\pi\epsilon_0\,r^2}\,\dot{r} = -\frac{e^6}{12\pi^3\epsilon_0^3\,m_e^2\,c^3\,r^4} \; .
\end{equation}

Rearranging:
\begin{equation}
  r^2\,\mathrm{d}r = -\frac{2e^4}{3\,(4\pi\epsilon_0)^2\,m_e^2\,c^3}\,\mathrm{d}t \;\equiv\; -C\,\mathrm{d}t \; .
  \label{eq:ode}
\end{equation}

% Subsection ...........................................................................................................
\subsection{Integration and Result}
\label{subsection:result}

Integrating~(\ref{eq:ode}) from $r_0 = a_0$ (the Bohr radius, $a_0 \approx 0.529\,\text{\AA}$) to $r = 0$:
\begin{equation}
  \int_{a_0}^{0} r^2\,\mathrm{d}r = -C\,t_{\mathrm{collapse}} \implies \frac{a_0^3}{3} = C\,t_{\mathrm{collapse}},
\end{equation}

giving the final results
\begin{equation}
  \boxed{t_{\mathrm{collapse}} = \frac{(4\pi\epsilon_0)^2\,m_e^2\,c^3\,a_0^3}{4\,e^4} \approx 1.6\times10^{-11}\,\mathrm{s} \; .}
  \label{eq:tcollapse}
\end{equation}

Sixteen picoseconds. This is not an order-of-magnitude estimate; it is the exact result of integrating three independently verified and internally consistent physical laws. The calculation is impeccable, yet the conclusion is empirically false. That tension is the starting point of everything that follows.

\begin{table}[h]
  \centering
  \begin{tabular}{|l|l|l|}
    \toprule
    Constant & Symbol & Value \\
    \midrule
    Bohr radius            & $a_0$        & $0.529\times10^{-10}$\,m \\
    Electron mass          & $m_e$        & $9.109\times10^{-31}$\,kg \\
    Elementary charge      & $e$          & $1.602\times10^{-19}$\,C \\
    Speed of light         & $c$          & $2.998\times10^{8}$\,m\,s$^{-1}$ \\
    Permittivity of vacuum & $\epsilon_0$ & $8.854\times10^{-12}$\,F\,m$^{-1}$ \\
    \bottomrule
  \end{tabular}
  \caption{Physical constants used in the collapse calculation.}
  \label{tab:constants}
\end{table}
